% ============================================================
% MÓDULO 2: CADENAS DE MARKOV
% ============================================================
\section{M2 · Cadenas de Markov}
\module{Módulo 2: Cadenas de Markov}{Modelando la evolución estocástica del entorno cuando no podemos intervenir}

\begin{frame}{OptiBot y la Cinta Transportadora}
  \begin{columns}[c,onlytextwidth]
    \column{0.50\textwidth}
      \textbf{El Escenario: Inspección de Rodillos}
      \begin{itemize}
        \item OptiBot vigila la cinta central de paquetes.
        \item Los rodamientos pasan por 3 estados:
        \begin{enumerate}
          \item \state{Nuevo ($N$)}
          \item \state{Desgastado ($D$)}
          \item \state{Averiado ($A$)}
        \end{enumerate}
        \item Al final de cada turno se inspeccionan.
        \item \textbf{Pregunta:} Si hoy está nuevo, ¿cuál es la probabilidad de que esté averiado en 2 turnos?
      \end{itemize}
    \column{0.48\textwidth}
      \centering
      \begin{tikzpicture}[scale=0.85,transform shape,
        st/.style={circle,draw=Ink,fill=Teal!12,thick,minimum size=1.1cm,font=\bfseries},
        edge/.style={-{Latex[length=2.5mm]},thick,SoftText}]
        \node[st] (N) at (0,2) {$N$};
        \node[st,fill=Gold!25] (D) at (2.5,2) {$D$};
        \node[st,fill=Coral!20] (A) at (1.25,0) {$A$};

        \draw[edge] (N) to[loop left] node[left] {$0.7$} (N);
        \draw[edge,Teal] (N) to[bend left=15] node[above] {$0.2$} (D);
        \draw[edge] (N) to[bend right=15] node[below left] {$0.1$} (A);

        \draw[edge] (D) to[loop right] node[right] {$0.6$} (D);
        \draw[edge,Coral] (D) to[bend left=15] node[below right] {$0.3$} (A);
        \draw[edge] (D) to[bend left=15] node[below] {$0.1$} (N);

        \draw[edge] (A) to[loop below] node[below] {$0.9$} (A);
        \draw[edge] (A) to[bend right=30] node[left] {$0.1$} (N);
      \end{tikzpicture}
  \end{columns}
  \vspace{0.3cm}
  \keyidea{La \textbf{Propiedad de Markov}: Dado el estado actual, el futuro es condicionalmente independiente del pasado.}
\end{frame}

\begin{frame}{La Matriz de Transición de Estados}
  \begin{columns}[c,onlytextwidth]
    \column{0.50\textwidth}
      \begin{block}{Definición Formal}
        Para cualquier secuencia de estados:
        \[
          P(S_{t+1} \mid S_t, S_{t-1}, \dots, S_0) = P(S_{t+1} \mid S_t)
        \]
      \end{block}
      \vspace{0.2cm}
      Agrupamos todas las probabilidades de transición en la matriz estocástica $P$:
      \[
        P_{ij} = P(S_{t+1} = j \mid S_t = i)
      \]
      \textbf{Regla de oro:} Cada fila debe sumar exactamente 1.
    \column{0.46\textwidth}
      \[
        P = \begin{pmatrix}
          0.7 & 0.2 & 0.1 \\
          0.1 & 0.6 & 0.3 \\
          0.1 & 0.0 & 0.9
        \end{pmatrix}
      \]
      \begin{itemize}
        \item Fila 1 ($N$): $0.7 + 0.2 + 0.1 = 1.0$
        \item Fila 2 ($D$): $0.1 + 0.6 + 0.3 = 1.0$
        \item Fila 3 ($A$): $0.1 + 0.0 + 0.9 = 1.0$
      \end{itemize}
      \vspace{0.2cm}
      \small ¿Probabilidad en 2 pasos? Simplemente calculamos $P^2 = P \times P$.
  \end{columns}
\end{frame}

\begin{frame}{Cálculo Manual: Distribución a 2 Pasos y Estado Estacionario}
  \begin{columns}[T,onlytextwidth]
    \column{0.48\textwidth}
      \begin{exampleblock}{Cálculo a 2 turnos desde $N$}
        Partiendo de $\mu_0 = [1, 0, 0]$:
        \begin{align*}
          \mu_1 &= \mu_0 P = [0.7, 0.2, 0.1] \\
          \mu_2 &= \mu_1 P \\
          P(A \text{ en } t=2) &= (0.7)(0.1) + (0.2)(0.3) + (0.1)(0.9) \\
          &= 0.07 + 0.06 + 0.09 = \mathbf{0.22}
        \end{align*}
      \end{exampleblock}
      \small Hay un $22\%$ de probabilidad de fallo al segundo turno.
    \column{0.48\textwidth}
      \begin{block}{Distribución Estacionaria $\pi^*$}
        A largo plazo, el sistema alcanza el equilibrio:
        \[
          \pi^* P = \pi^* \quad \text{con} \quad \sum_i \pi^*_i = 1
        \]
      \end{block}
      \vspace{0.2cm}
      \small Nos dice qué porcentaje del tiempo la cinta pasará averiada, sin importar dónde comenzó.
  \end{columns}
  \vspace{0.35cm}
  \keyidea{El estado estacionario describe la distribución de probabilidad invariante bajo la dinámica del sistema.}
\end{frame}

\begin{frame}[fragile]{Verificación Numérica en Python}
  \begin{columns}[T,onlytextwidth]
    \column{0.62\textwidth}
\begin{lstlisting}[style=rlcode]
import numpy as np

P = np.array([
    [0.7, 0.2, 0.1],
    [0.1, 0.6, 0.3],
    [0.1, 0.0, 0.9]
])

# Proyeccion a 2 turnos desde Nuevo
mu_0 = np.array([1.0, 0.0, 0.0])
mu_2 = mu_0 @ np.linalg.matrix_power(P, 2)
print("Estado en t=2:", np.round(mu_2, 3))
# [0.52 0.26 0.22] -> 22% de averia confirmada

# Estado estacionario: autovector con autovalor 1
evals, evecs = np.linalg.eig(P.T)
stationary = evecs[:, np.isclose(evals, 1.0)].real
stationary = stationary / np.sum(stationary)
print("Distribucion estacionaria pi*:", np.round(stationary.T[0], 3))
# [0.25 0.125 0.625]
\end{lstlisting}
    \column{0.35\textwidth}
      \begin{alertblock}{La limitación de M2}
        Ahora sabemos cómo evoluciona el mundo en el tiempo y podemos predecir el futuro.\\
        \vspace{0.2cm}
        \textbf{Pero somos espectadores pasivos}: OptiBot sólo mira la cinta romperse. No puede intervenir con acciones.\\
        \vspace{0.2cm}
        ¿Y si unimos las \textbf{acciones} de M1 con los \textbf{estados} de M2? $\to$ \textbf{¡MDP!}
      \end{alertblock}
  \end{columns}
\end{frame}
